Class imbalance makes credit assessment and thyroid screening difficult to evaluate using accuracy alone. We compared decision trees, random forests, logistic regression and support vector classifiers on two UCI datasets. Each model was tested with plain training, SMOTE, class weighting and threshold adjustment. We used the same validation splits and fitted preprocessing within training folds. Apriori supplied association rules to help explain patterns in German Credit. Cost sensitive logistic regression achieved the highest mean F1 in the German Credit main experiment, while random forest with SMOTE led on sick. Strategies that improved recall could reduce precision, showing why accuracy alone can hide poor minority class performance. The additional ensemble study achieved a higher exploratory F1 on sick, but we did not directly test its advantage over the main experiment winner. Apriori found combinations associated with higher bad proportions and added to the tree model explanations. These results show that model choice depends on the dataset and decision objective. Independent validation is needed before practical use.
